\section{問題定義}\label{sec:problem}

本節では、密集区間予測に必要な概念を形式的に定義する。

\subsection{密集区間}

\begin{definition}[密集区間]
  トランザクションデータベース $\mathcal{D} = \{T_0, T_1, \ldots, T_{n-1}\}$、
  アイテムセット $X$、ウィンドウサイズ $W$、閾値 $\theta$ が与えられたとき、
  区間 $[s, e]$ がアイテムセット $X$ の\textbf{密集区間}であるとは、
  全ての $l \in [s, e]$ に対して
  \begin{equation}
    |\{i \mid l \leq i \leq l + W,\; X \subseteq T_i\}| \geq \theta
  \end{equation}
  が成り立ち、かつ $[s, e]$ が極大であることを言う。
  実際の密集期間は $[s, e + W]$ である。
\end{definition}

\subsection{密集区間発生過程（DIOP）}

\begin{definition}[密集区間発生過程]
  アイテムセット $X$ に対する密集区間の列
  $\mathcal{I}_X = \{(s_1, e_1), (s_2, e_2), \ldots, (s_K, e_K)\}$
  （$s_1 < s_2 < \cdots < s_K$）が与えられたとき、
  到着時刻列 $\{s_1, s_2, \ldots, s_K\}$ を
  \textbf{密集区間発生過程}（Dense Interval Occurrence Process; DIOP）と呼ぶ。
\end{definition}

DIOPはカウント過程 $N(t) = |\{k \mid s_k \leq t\}|$ として表現でき、
条件付き強度関数 $\lambda(t|\mathcal{H}_t)$ によって特徴づけられる。
ここで $\mathcal{H}_t$ は時刻 $t$ までの履歴である。

\subsection{密集区間間時間（IDIT）}

\begin{definition}[密集区間間時間]
  密集区間列 $\mathcal{I}_X$ に対し、第 $k$ 番目の
  \textbf{密集区間間時間}（Inter-Dense Interval Time; IDIT）を
  \begin{equation}
    \tau_k = s_{k+1} - (e_k + W), \quad k = 1, \ldots, K-1
  \end{equation}
  と定義する。
  ここで $e_k + W$ は第 $k$ 番目の密集区間の実質的な終了時刻である。
  $\tau_k < 0$ の場合（密集区間が重複する場合）は $\tau_k = 0$ とする。
\end{definition}

IDITの分布は密集区間の発生パターンを特徴づける重要な統計量である。
変動係数 $\mathrm{CV} = \sigma_\tau / \mu_\tau$ は発生パターンの規則性を示す：
\begin{itemize}
  \item $\mathrm{CV} \approx 0$：密集区間は等間隔に近い（規則的）
  \item $\mathrm{CV} \approx 1$：Poisson的（ランダム）
  \item $\mathrm{CV} > 1$：クラスタリングが存在（バースト的）
\end{itemize}

\subsection{密集持続時間}

\begin{definition}[密集持続時間]
  密集区間 $(s_k, e_k)$ の\textbf{密集持続時間}を
  \begin{equation}
    d_k = e_k - s_k + W
  \end{equation}
  と定義する。$d_k$ は密集条件が維持される実質的な期間長である。
\end{definition}

\subsection{予測問題の定式化}

本研究では以下の2つの予測問題を扱う。

\begin{problem}[次回発生時刻予測]\label{prob:arrival}
  密集区間の履歴 $\mathcal{I}_X^{(1:K)} = \{(s_1, e_1), \ldots, (s_K, e_K)\}$ が
  与えられたとき、次の密集区間の開始時刻 $s_{K+1}$ を予測せよ。
  具体的には、条件付き強度関数 $\lambda(t|\mathcal{H}_{s_K})$ のもとでの
  次回到着時刻の条件付き期待値
  $\hat{s}_{K+1} = \mathbb{E}[s_{K+1} | \mathcal{H}_{s_K}]$
  を求める。
  予測の精度は、ベースライン手法（定常Poisson過程、IDIT平均値外挿）との比較で評価する。
\end{problem}

\begin{problem}[持続時間予測]\label{prob:duration}
  密集区間の持続時間の履歴 $\{d_1, d_2, \ldots, d_K\}$ と
  オプションとして共変量ベクトル $\mathbf{x}_k$（前回のIDIT、前回の持続時間等）が
  与えられたとき、次の密集区間の持続時間 $d_{K+1}$ を予測せよ。
  予測の精度はMAE、RMSEで評価し、ナイーブベースライン（標本平均）と比較する。
\end{problem}

\subsection{評価指標}

予測問題の評価には以下の指標を用いる。
\begin{itemize}
  \item \textbf{相対誤差}（RE）：$\mathrm{RE} = |{\hat{s}_{K+1} - s_{K+1}}| / {s_{K+1}}$
  \item \textbf{平均絶対誤差}（MAE）：$\mathrm{MAE} = \frac{1}{n}\sum_{i=1}^{n}|\hat{d}_i - d_i|$
  \item \textbf{対数尤度比}（LLR）：$\mathrm{LLR} = \ell_{\text{Hawkes}} - \ell_{\text{Poisson}}$。
    自己励起構造の有無に対する尤度比検定に使用。$\mathrm{LLR} > 0$ かつ
    カイ二乗検定で有意（$p < 0.05$）の場合にHawkesモデルが優位と判定する。
\end{itemize}
